% ======================================================================
% CHAPTER 2: LITERATURE REVIEW AND THEORETICAL FOUNDATIONS
% ======================================================================

\chapter{Literature Review and Theoretical Foundations}
\label{chap:literature_review}

\section{Atmospheric Hydrodynamics and Dust Saltation Physics}
\label{sec:physics_foundations}

Sand and dust storm events represent complex, multi-scale physical phenomena governed by the coupled dynamics of planetary boundary layer fluid mechanics, surface aeolian erosion, and multi-phase aerosol microphysics. Establishing high-accuracy medium- to long-term forecasting models requires a rigorous mathematical understanding of the first-principles laws controlling particle detachment, turbulent vertical lofting, long-range advective dispersion, and dry/wet deposition.

\subsection{Boundary Layer Turbulence and Rotating Navier-Stokes Equations}
\label{subsec:navier_stokes}

Large-scale atmospheric circulation is governed by the conservation of momentum in a rotating planetary reference frame, described by the Navier-Stokes equations~\cite{shao2008physics}:
\begin{equation}
\label{eq:navier_stokes}
\frac{\partial \mathbf{u}}{\partial t} + (\mathbf{u} \cdot \nabla)\mathbf{u} = -\frac{1}{\rho_{\text{air}}}\nabla p + \mathbf{g} - 2\mathbf{\Omega} \times \mathbf{u} + \nu \nabla^2 \mathbf{u}
\end{equation}
where $\mathbf{u} = (u, v, w)^T$ is the three-dimensional wind velocity vector, $\rho_{\text{air}}$ is moist air density, $p$ is static atmospheric pressure, $\mathbf{g} = (0, 0, -g)^T$ represents the gravitational acceleration vector, $\mathbf{\Omega}$ denotes the planetary angular rotation vector (inducing the Coriolis acceleration term $-2\mathbf{\Omega} \times \mathbf{u}$), and $\nu$ is kinematic molecular viscosity.

In the Atmospheric Boundary Layer (ABL, the lowest $1\text{--}2\,\text{km}$ of the troposphere), molecular viscous stresses are entirely dominated by turbulent Reynolds stresses. The surface shear stress $\tau_0$ exerted by boundary layer air streams on ground surfaces provides the physical kinetic energy required to dislodge mineral particles. In micrometeorology, the surface friction velocity $u_*$ is formally defined as:
\begin{equation}
\label{eq:friction_velocity_def}
u_* = \sqrt{\frac{\tau_0}{\rho_{\text{air}}}}
\end{equation}
According to Monin-Obukhov Similarity Theory (MOST), within a horizontally homogeneous surface constant-flux layer, the vertical profile of mean horizontal wind speed $U(z)$ as a function of height $z$ satisfies the logarithmic wind law:
\begin{equation}
\label{eq:monin_obukhov}
U(z) = \frac{u_*}{\kappa} \left[ \ln\left(\frac{z}{z_0}\right) - \psi_m\left(\frac{z}{L_{\text{MO}}}\right) \right]
\end{equation}
where $\kappa \approx 0.40$ is the dimensionless von K\'arm\'an constant, $z_0$ is the aerodynamic surface roughness length, $L_{\text{MO}}$ denotes the Monin-Obukhov stability length scale, and $\psi_m$ is the empirical universal thermal stability correction function ($\psi_m > 0$ under unstable convective conditions, $\psi_m = 0$ for neutral conditions, and $\psi_m < 0$ under stable inversion layers). During spring over East Asian desert margins, intense diurnal solar insolation heats the denuded ground, generating superadiabatic lapse rates and driving positive $\psi_m$ excursions that violently transfer downward momentum to the ground.

\subsection{Aerodynamic Saltation Mechanics: Threshold Shear Stress and Kinetic Bombardment}
\label{subsec:saltation_mechanics}

Foundational aeolian physics established by Bagnold~\cite{bagnold1941physics} and Owen~\cite{owen1964saltation} revealed that ultra-fine clay and silt particles ($d_p < 20\,\mu\text{m}$) cannot be lifted directly from smooth surfaces by aerodynamic drag alone, due to strong inter-particle Van der Waals cohesion. Instead, dust emission is driven primarily by **saltation** (grain bombardment): coarse sand grains ($d_p \approx 70\text{--}500\,\mu\text{m}$) are mobilized first, hopping along ballistic trajectories within the near-surface saltation layer ($0.1\text{--}1\,\text{m}$ height). As these energetic sand grains impact the ground, their kinetic energy fractures aggregate bonds, splashing fine sub-micron dust particles into turbulent boundary layer eddies.

Saltation initiation requires the surface friction velocity $u_*$ to exceed a critical aerodynamic threshold friction velocity $u_{*t}$. Based on force balances between aerodynamic drag, lift, gravity, and cohesive inter-particle forces formulated by Shao~\cite{shao2008physics}, Gillette~\cite{gillette1979environmental}, and Kok~\cite{kok2012physics}, $u_{*t}$ is expressed as:
\begin{equation}
\label{eq:threshold_friction_velocity}
u_{*t} = A_N \sqrt{\frac{\rho_{\text{sand}} - \rho_{\text{air}}}{\rho_{\text{air}}} g d_p + \frac{\gamma}{\rho_{\text{air}} d_p}} \cdot f(w_s) \cdot f(z_0)
\end{equation}
where $A_N \approx 0.11$ is a dimensionless empirical scaling coefficient, $\rho_{\text{sand}} \approx 2{,}650\,\text{kg/m}^3$ is quartz particle density, and $\gamma \approx 1.65 \times 10^{-4}\,\text{N/m}$ is the surface inter-particle cohesion coefficient.

The function $f(w_s)$ accounts for the capillary binding effect of soil moisture in resisting aeolian detachment~\cite{ref13}:
\begin{equation}
\label{eq:moisture_inhibition}
f(w_s) = 
\begin{cases}
1, & w_s \le w_s' \\
\sqrt{1 + 1.21 \left(w_s - w_s'\right)^{0.68}}, & w_s > w_s'
\end{cases}
\end{equation}
where $w_s$ is volumetric topsoil water content ($\text{m}^3/\text{m}^3$) and $w_s'$ is the threshold bound-water moisture capacity. When topsoil is damp, meniscus capillary tension bonds soil grains, causing threshold friction velocities to increase exponentially.

Once $u_* > u_{*t}$, the horizontal saltation mass flux $F_{\text{salt}}$ ($\text{kg}/(\text{m}\cdot\text{s})$) scales with the cube of friction velocity, following Owen's (1964) analytical closed-form solution~\cite{owen1964saltation}:
\begin{equation}
\label{eq:owen_flux}
F_{\text{salt}} = c_{\text{salt}} \frac{\rho_{\text{air}}}{g} u_*^3 \left(1 - \frac{u_{*t}^2}{u_*^2}\right) \cdot \mathbb{I}(u_* > u_{*t})
\end{equation}
where $c_{\text{salt}} \approx 1.6\text{--}2.3$ is the Owen coefficient and $\mathbb{I}(\cdot)$ is the Heaviside indicator function. The vertical dust emission flux $F_{\text{dust}}$ ($\mu\text{g}/(\text{m}^2\cdot\text{s})$) injected into the atmosphere is proportional to horizontal saltation via the sandblasting efficiency factor $\alpha_{\text{sand}}$:
\begin{equation}
\label{eq:vertical_emission_flux}
F_{\text{dust}} = \alpha_{\text{sand}} \cdot F_{\text{salt}} = \alpha_{\text{sand}} \cdot c_{\text{salt}} \frac{\rho_{\text{air}}}{g} u_*^3 \left(1 - \frac{u_{*t}^2}{u_*^2}\right) \cdot \mathbb{I}(u_* > u_{*t})
\end{equation}
Equation~(\ref{eq:vertical_emission_flux}) represents the core aerodynamic gating law of wind erosion: **dust emission is strictly zero when $u_* \le u_{*t}$**, exploding cubically with friction velocity once threshold shear is breached.

\subsection{Three-Dimensional Aerosol Advection-Diffusion-Deposition Governing Equation}
\label{subsec:aerosol_pde}

Following turbulent vertical injection, suspended mineral dust particulates are transported through the troposphere. In an Eulerian frame of reference, the conservation of airborne particulate mass concentration $C(\mathbf{x}, t)$ ($\mu\text{g/m}^3$) is governed by the three-dimensional advection-diffusion-deposition partial differential equation (PDE):
\begin{equation}
\label{eq:aerosol_full_pde}
\underbrace{\frac{\partial C}{\partial t}}_{\text{Local Tendency}} + \underbrace{\nabla \cdot (\mathbf{u} C)}_{\text{Advective Flux}} = \underbrace{\nabla \cdot (\mathbf{K} \nabla C)}_{\text{Turbulent Diffusion}} + \underbrace{S_{\text{emission}}(\mathbf{x}, t)}_{\text{Surface Source}} - \underbrace{v_d C}_{\text{Dry Deposition}} - \underbrace{\Lambda_{\text{scav}} P_{\text{rain}} C}_{\text{Wet Scavenging}}
\end{equation}
where:
\begin{itemize}
    \item $\mathbf{u} = (u, v, w)^T$ is the three-dimensional wind advection vector;
    \item $\mathbf{K} = \text{diag}(K_{xx}, K_{yy}, K_{zz})$ represents the turbulent eddy diffusivity tensor, with vertical mixing $K_{zz}$ parameterized using boundary layer turbulent closure schemes (e.g., Mellor-Yamada);
    \item $S_{\text{emission}}$ represents the surface saltation flux defined in Equation~(\ref{eq:vertical_emission_flux}), active exclusively at the surface boundary;
    \item $v_d$ denotes Stokes gravitational terminal settling velocity, modulated by particle diameter $d_p$, particle density, and the Cunningham slip correction factor;
    \item $\Lambda_{\text{scav}}$ is the empirical scavenging coefficient and $P_{\text{rain}}$ is precipitation rate ($\text{mm/h}$).
\end{itemize}

Equation~(\ref{eq:aerosol_full_pde}) provides the physical mass-conservation governing constraint embedded within the Physics-Informed Neural Network (PINN) architecture developed in subsequent chapters.

---

\section{Numerical Weather Prediction and Aerosol Modeling Systems}
\label{sec:nwp_modeling_review}

\subsection{Operational NWP and Regional Chemical Transport Models}
\label{subsec:nwp_systems}

For decades, Eulerian numerical models have served as the foundation of operational atmospheric dust forecasting~\cite{ref2}. State-of-the-art numerical modeling suites include:
\begin{enumerate}
    \item \textbf{ECMWF Integrated Forecasting System (IFS-AER)}: A global numerical weather prediction suite coupling dynamic aerosol assimilation of MODIS and Sentinel Aerosol Optical Depth (AOD) at $0.1^\circ$ resolution ($\approx 9\,\text{km}$).
    \item \textbf{CMA-GFS and CUACE/Dust (China Meteorological Administration)}: An operational model developed by the Chinese Academy of Meteorological Sciences incorporating East Asian desert geographical databases and dynamic wind erosion formulations~\cite{ref2,zhang2004evolution}.
    \item \textbf{WRF-Chem (Weather Research and Forecasting with Chemistry)}: A non-hydrostatic regional model supporting online aerosol-cloud-radiation feedbacks, widely employed for high-resolution case analyses.
\end{enumerate}

\subsection{Dust Emission Parameterization Schemes}
\label{subsec:emission_schemes}

Within numerical models, dust source terms rely on parameterization schemes. Table~\ref{tab:dust_schemes_comparison} presents a comparative synthesis of four widely adopted dust emission parameterization schemes.

\begin{table}[H]
\centering
\small
\caption{Comparative Analysis of Mainstream Dust Emission Parameterization Schemes}
\label{tab:dust_schemes_comparison}
\begin{tabularx}{\textwidth}{l p{2.8cm} p{3.2cm} p{3.5cm} X}
\toprule
\textbf{Scheme} & \textbf{Key Reference} & \textbf{Core Physical Hypothesis} & \textbf{Input Parameters} & \textbf{Operational Limitations} \\
\midrule
Gillette Scheme & Gillette (1979)~\cite{gillette1979environmental} & Empirical vertical-to-horizontal flux ratio & $u_*, u_{*t}$, soil clay fraction & Neglects kinetic particle energy dissipation; overestimates sand flux \\
Ginoux Scheme & Ginoux et al. (2001) & Source erodibility matrix with quartic wind scaling & $U_{10}$, fractional vegetation cover, soil classifications & Lacks dynamic surface roughness adjustment; lagged spring thaw response \\
Shao Scheme & Shao (2008)~\cite{shao2008physics} & Saltation impact energy fracturing cohesive aggregate bonds & Size-resolved particle distribution, plastic bond energy, moisture & High parameter calibration complexity; sensitive to sub-grid wind bursts \\
Kok Scheme & Kok (2012)~\cite{kok2012physics} & Fracture mechanics and non-elastic energy dissipation & Soil aggregate brittleness, $u_*, u_{*t}$, binding energy & Requires microstructural soil texture data; highly sensitive to gustiness \\
\bottomrule
\end{tabularx}
\end{table}

As shown in Table~\ref{tab:dust_schemes_comparison}, empirical emission parameterizations rely on stationary wind-tunnel calibrations. When applied across heterogeneous desert landscapes, small errors in input soil parameters propagate non-linearly, leading to severe flux overestimation or underestimation.

\subsection{Extended-Range Forecasts, Atmospheric Chaos, and the Lyapunov Horizon}
\label{subsec:lyapunov_horizon}

In dynamical meteorology, non-linear chaos imposes an absolute theoretical ceiling on deterministic forecasting skill. As demonstrated by Lorenz (1963)~\cite{lorenz1963deterministic}, atmospheric error growth in phase space is governed by the leading Lyapunov exponent ($\lambda_{\text{max}}$):
\begin{equation}
\label{eq:lyapunov_divergence}
\|\delta \mathbf{x}(t)\| \approx \|\delta \mathbf{x}(0)\| \cdot \exp(\lambda_{\text{max}} t)
\end{equation}
For synoptic-scale baroclinic waves in mid-latitudes, the e-folding error doubling time ($1/\lambda_{\text{max}}$) is approximately $5\text{--}7\,\text{days}$. The predictable deterministic horizon $T_{\text{Lyap}}$ for fine particulate transport under active cyclogenesis collapses to:
\begin{equation}
\label{eq:t_lyap}
T_{\text{Lyap}} \approx 72\,\text{hours} \quad (3\,\text{days})
\end{equation}
Beyond 72 hours, deterministic NWP wind trajectories decouple from physical reality, causing modeled aerosol plumes to disperse prematurely or miss narrow transport corridors entirely. Consequently, extending skillful forecasts to the 3--15 day window requires integrating statistical residual post-processing (Main Line A) with deep spatiotemporal foundation priors (Main Line B).

---

\section{Machine Learning and Deep Spatiotemporal Neural Networks in Atmospheric Forecasting}
\label{sec:ml_review}

The rapid accumulation of petabyte-scale earth observations has established data-driven machine learning as a transformative pillar of modern meteorology~\cite{ref1,ref32,ref40}. This section reviews the chronological and architectural evolution of AI methodologies in atmospheric and sandstorm forecasting.

\subsection{Historical Evolution: Early Shallow Learners in Sandstorm Identification}
\label{subsec:early_ann_svm}

Initial applications of machine learning focused on statistical point-forecasting at isolated meteorological stations. Huang et al.~\cite{ref3} pioneered an Artificial Neural Network (ANN) employing stepwise multiple linear regression to screen pressure, temperature, humidity, and wind speed features, achieving superior skill over conventional linear regressions in northwestern China. Lu et al.~\cite{ref4} introduced Support Vector Machines (SVM) driven by 500\,hPa geopotential height and low-level wind fields to classify sandstorm days.

However, these early investigations identified a fundamental hurdle: sandstorms represent extreme, rare events within multi-year records. Normal days vastly outnumber dust days, yielding severe **class imbalance**. To prevent models from collapsing into trivial non-event predictions, Xie et al.~\cite{ref5} proposed an adaptive hybrid resampling strategy, while Zhang et al.~\cite{ref6} integrated Synthetic Minority Over-sampling Technique (SMOTE) with AdaBoost random forest ensembles, markedly improving detection rates for rare dust events.

Subsequent investigations benchmarked diverse algorithm families. Kaboodvandpour et al.~\cite{ref7} evaluated multiple regressions, Adaptive Neuro-Fuzzy Inference Systems (ANFIS), and random forests across stations in western Iran. Shaiba et al.~\cite{ref9} and Al Murayziq et al.~\cite{ref8} applied Bayesian belief networks and Case-Based Reasoning (CBR) to classify imminent dust storms. Ebrahimi-Khusfi et al.~\cite{ref13} examined multi-factor drivers—including Normalized Difference Vegetation Index (NDVI)~\cite{ref14} and soil characteristics—to forecast monthly dust storm frequencies. Mahmoudi et al.~\cite{ref15} combined 29 years of station observations with Radial Basis Function (RBF) networks and TOPSIS multi-criteria decision modeling to map spatial dust hazard vulnerabilities in ArcGIS.

\subsection{Spatiotemporal Sequence Deep Learning and Atmospheric Foundation Models}
\label{subsec:deep_spatiotemporal}

Over the past decade, deep architectures capable of modeling complex spatiotemporal dynamics have transformed meteorological forecasting:
\begin{itemize}
    \item \textbf{Spatiotemporal Recurrent Networks}: Shi et al.~\cite{ref17} introduced Convolutional LSTM (ConvLSTM), embedding convolutional operations within LSTM recurrent memory cells to capture spatiotemporal precipitation dynamics. Wang et al.~\cite{ref18} advanced this paradigm with PredRNN, developing Spatiotemporal Memory Cells (ST-LSTM) that decouple spatial memory from temporal state transitions.
    \item \textbf{Attention and Spatiotemporal Transformers}: To overcome the sequential processing bottlenecks of recurrent networks, Bai et al.~\cite{ref19} proposed Rainformer for radar-based precipitation forecasting. Gao et al.~\cite{ref20} introduced Earthformer, introducing cuboid self-attention to model multi-scale earth system interactions. In tropical cyclone tracking, Huang et al.~\cite{ref21,ref22} developed MMSTN and MGTCF to handle heterogeneous multi-modal meteorological observations.
    \item \textbf{Deep Learning in Sandstorm Modeling}: Li et al.~\cite{ref10} proposed the INB-CNN model, fusing naive Bayes (for surface features) with CNNs (for circulation fields). Harba et al.~\cite{ref11} utilized Feature Pyramid Networks (FPN) and bounding-box detection to track dust plume boundaries from satellite imagery. Ren et al.~\cite{ref12} applied transfer learning from pre-trained image weights to accelerate CNN convergence for Inner Mongolian sandstorms. Ren et al.~\cite{ref16} introduced a Stacking ensemble fusing LSTM and CNN backbones, demonstrating marked improvements over standalone baselines.
\end{itemize}

\subsection{Non-Euclidean Graph Neural Networks for Atmospheric Transport Corridors}
\label{subsec:gnn_review}

While convolutional architectures excel on regular 2D pixel grids, continental atmospheric dust transport is intrinsically non-Euclidean. Narrow orographic passes—such as the Hexi Corridor, the Helan Mountain pass, and the Yan-Taihang gap—act as physical funnels, steering intense dust fluxes along discrete topological pathways. Standard grid convolutions treat space as homogeneous and flat, which diffuses sharp plume boundaries and introduces unnecessary computational overhead.

Building on Graph Convolutional Networks (GCN) developed by Kipf and Welling~\cite{kipf2017semi}, Graph Neural Networks (GNNs) perform spectral filtering and message-passing over arbitrary graph topologies $\mathcal{G} = (\mathcal{V}, \mathcal{E})$. By representing dust sources, transport funnels, and receptor metropolitan clusters as nodes connected by dynamic, wind-directed edges, GNNs model orographic channeling and advective momentum with high spatial fidelity and computational efficiency.

\subsection{Cross-Domain Applications of Machine Learning in Weather-Dependent Sectors}
\label{subsec:cross_domain_ml}

Machine learning methodologies have matured across diverse weather-dependent domains, providing valuable algorithmic templates for multi-source environmental modeling:
\begin{enumerate}
    \item \textbf{Severe Convective Weather, Gales, and Fog}: Automated Machine Learning (AutoML) frameworks (such as AutoGluon) achieved 96.7\% hit rates in predicting thunderstorm gales~\cite{ref25,ref26}. The hybrid AutoFeat-TPOT algorithm delivered accurate short-term fog classification~\cite{ref27}. Tao et al.~\cite{ref28} and Zhao et al.~\cite{ref29} utilized random forests and the AeolusStorm system to forecast convective shear at aerospace launch facilities. Luo and Duan~\cite{ref30} applied firefly-optimized SVMs for multi-scale weather hazard mapping. Pu et al.~\cite{ref31} and Guan and Su~\cite{ref32} integrated decision trees with numerical models for highway icing and multi-hazard forecasting. Zhao et al.~\cite{ref33} combined multi-layer perceptrons with blockchain logging for flood forecasting, while Zhang et al.~\cite{ref34} deployed CART trees to predict snowfall accumulation.
    \item \textbf{Atmospheric Parameters and Boundary Layer Diagnostics}: Bai et al.~\cite{ref35} fused ANN, SVM, and LSTM models to predict Zenith Tropospheric Delay (ZTD) for satellite navigation. Bai et al.~\cite{ref36} applied XGBoost to continuous planetary boundary layer height (PBLH) retrieval using wind profiler telemetry. An et al.~\cite{ref37} benchmarked solar irradiance prediction algorithms. Liu~\cite{ref38} formulated the EALSTM-QR model for wind speed prediction intervals, and Mahavik et al.~\cite{ref39} utilized random forests for radar precipitation bias correction. Google's NeuralGCM~\cite{ref40} demonstrated that hybrid machine learning can match or exceed traditional general circulation models in computational efficiency and medium-range skill.
    \item \textbf{Cross-Industry Operations (Energy, Transportation, and Agriculture)}: In energy management, CatBoost, LightGBM, and LSTM architectures are widely applied for electrical building load forecasting~\cite{ref41}, extreme-weather Bayesian load modeling~\cite{ref42}, ultra-short-term wind and photovoltaic power generation~\cite{ref43,ref44,ref45}, district heating temperature control~\cite{ref46,ref47,ref48}, and natural gas demand forecasting~\cite{ref49}. In transportation, graph convolutions and Transformers predict urban rail transit passenger surges~\cite{ref50}, while SHAP-interpreted XGBoost models analyze expressway interchange crash risks~\cite{ref51} and maritime safety~\cite{ref52}. In precision agriculture, ensemble algorithms (such as RFXG) utilize weather indices for crop yield forecasting and nutrient recommendations across cotton, sugarcane, and rice cropping systems~\cite{ref53,ref54,ref55,ref56,ref57,ref58,ref59}. Finally, machine learning techniques underpin operational $\text{PM}_{2.5}$ concentration warnings~\cite{ref1,ref60} and industrial volatile organic compound (VOC) anomaly detection~\cite{ref61}.
\end{enumerate}

Despite these achievements, contemporary sandstorm research remains largely confined to short-range, local point classifications, lacking unified frameworks capable of continuous spatiotemporal modeling across the 3--15 day extended-range horizon.

---

\section{Physics-Informed Neural Networks (PINN) in Environmental Engineering}
\label{sec:pinn_review}

\subsection{Mathematical Foundations and Automatic Differentiation Residual Loss}
\label{subsec:pinn_architecture}

Purely data-driven neural networks minimize empirical risk over finite training datasets. When applied to rare extreme environmental events, unconstrained networks often converge to physically spurious local minima. Raissi et al.~\cite{raissi2019physics} formulated Physics-Informed Neural Networks (PINNs) to constrain neural representations using fundamental conservation laws.

Consider a general non-linear dynamical partial differential equation:
\begin{equation}
\label{eq:pinn_general_pde}
\mathcal{N}_\mathbf{x}[u(\mathbf{x}, t)] + \frac{\partial u(\mathbf{x}, t)}{\partial t} = 0, \quad \mathbf{x} \in \Omega, \quad t \in [0, T]
\end{equation}
where $u(\mathbf{x}, t)$ represents the physical state field (e.g., aerosol concentration or wind velocity) and $\mathcal{N}_\mathbf{x}[\cdot]$ is a non-linear differential operator containing spatial gradients, advection, and diffusion.

In a PINN, the solution is parameterized by a deep neural network $\hat{u}(\mathbf{x}, t; \theta) \approx u(\mathbf{x}, t)$ parameterized by weights $\theta$. Leveraging reverse-mode automatic differentiation (AD), the spatial and temporal partial derivatives (e.g., $\frac{\partial \hat{u}}{\partial t}, \nabla \hat{u}, \nabla^2 \hat{u}$) are computed analytically across the computational graph without discretization grids.

The physical differential residual $r(\mathbf{x}, t; \theta)$ is defined as:
\begin{equation}
\label{eq:pinn_residual}
r(\mathbf{x}, t; \theta) \coloneqq \frac{\partial \hat{u}(\mathbf{x}, t; \theta)}{\partial t} + \mathcal{N}_\mathbf{x}\left[\hat{u}(\mathbf{x}, t; \theta)\right]
\end{equation}
The composite PINN loss function is formulated as:
\begin{equation}
\label{eq:pinn_loss_formulation}
\mathcal{L}_{\text{PINN}}(\theta) = \mathcal{L}_{\text{data}}(\theta) + \lambda_{\text{phys}} \mathcal{L}_{\text{physics}}(\theta)
\end{equation}
where the empirical data loss is computed across sparse ground stations $\mathcal{D}_{\text{data}} = \{(\mathbf{x}_i, t_i, y_i)\}_{i=1}^{N_d}$:
\begin{equation}
\mathcal{L}_{\text{data}}(\theta) = \frac{1}{N_d} \sum_{i=1}^{N_d} \left\| \hat{u}(\mathbf{x}_i, t_i; \theta) - y_i \right\|^2
\end{equation}
and the physics loss is evaluated over arbitrary collocation points $\mathcal{D}_{\text{colloc}} = \{(\mathbf{x}_j^c, t_j^c)\}_{j=1}^{N_c}$ sampled uniformly across space and time:
\begin{equation}
\label{eq:pinn_colloc_loss}
\mathcal{L}_{\text{physics}}(\theta) = \frac{1}{N_c} \sum_{j=1}^{N_c} \left\| r(\mathbf{x}_j^c, t_j^c; \theta) \right\|^2
\end{equation}

Critically, the collocation set $\mathcal{D}_{\text{colloc}}$ **requires no observational ground truth labels**. Consequently, within unmonitored desert expanses, the physical residual forces the neural network to satisfy fundamental conservation laws at all virtual collocation coordinates.

\subsection{Resolving Dynamic Gradient Pathologies via Adaptive Balancing (GradNorm)}
\label{subsec:pinn_challenges}

Applying PINNs to continental sandstorm forecasting presents two technical challenges:
\begin{enumerate}
    \item \textbf{Differentiability of the Saltation Threshold Gate}: Aerodynamic sand detachment is governed by the non-linear saltation criterion $\mathbb{I}(u_* > u_{*t})$. The Heaviside step function possesses a discontinuous zero-gradient landscape everywhere except at the threshold, where its derivative is a Dirac delta function, which disrupts backpropagation. To resolve this, a smooth, continuous sigmoid approximation must replace the hard step function:
    \begin{equation}
    \sigma_\beta(u_*, u_{*t}) = \frac{1}{1 + \exp\left(-\beta (u_* - u_{*t})\right)}
    \end{equation}
    where $\beta$ controls transition steepness, providing well-conditioned gradients across the saltation boundary.
    \item \textbf{Gradient Pathologies in Multi-Objective Optimization}: During training, gradients originating from the data loss $\nabla_\theta \mathcal{L}_{\text{data}}$ and the differential PDE residual $\nabla_\theta \mathcal{L}_{\text{physics}}$ can diverge by several orders of magnitude. Static penalty multipliers $\lambda_{\text{phys}}$ frequently cause gradient vanishing or pathological oscillations. Dynamic gradient balancing algorithms (such as GradNorm) adaptively rescale multi-task loss weights, ensuring stable convergence toward physically consistent solutions.
\end{enumerate}

---

\section{Spatial Econometrics and Disaster Risk Behavioral Decision Theory}
\label{sec:spatial_behavioral}

\subsection{Spatial Autocorrelation Theory, Ordinary Kriging, and LISA Cluster Analysis}
\label{subsec:spatial_autocorrelation}

Atmospheric dust concentrations exhibit strong spatial continuity and regional spatial dependency. According to Tobler's First Law of Geography~\cite{tobler1970cellular}: ``All things are related, but near things are more related than distant things.''

Quantifying this spatial structure requires rigorous spatial econometric techniques:
\begin{enumerate}
    \item \textbf{Ordinary Kriging Geostatistical Interpolation}: Evaluates spatial covariance using an empirical semivariogram $\gamma(h)$:
    \begin{equation}
    \label{eq:semivariogram}
    \gamma(h) = \frac{1}{2 N(h)} \sum_{i=1}^{N(h)} \left[ Z(\mathbf{x}_i) - Z(\mathbf{x}_i + h) \right]^2
    \end{equation}
    Minimizing estimation variance under unbiased constraints produces continuous exposure surfaces from sparse monitoring stations.
    \item \textbf{Global Moran's $I$}: Quantifies overall spatial clustering across the study domain:
    \begin{equation}
    \label{eq:global_morans_i}
    I = \frac{n \sum_{i=1}^n \sum_{j=1}^n w_{ij} (z_i - \bar{z})(z_j - \bar{z})}{\left(\sum_{i=1}^n \sum_{j=1}^n w_{ij}\right) \sum_{i=1}^n (z_i - \bar{z})^2}
    \end{equation}
    where $w_{ij}$ represents the spatial weight matrix. A positive Moran's $I$ with $p < 0.01$ confirms statistically significant positive spatial clustering of dust hazards.
    \item \textbf{Local Indicators of Spatial Association (LISA)}: Decomposes global autocorrelation to identify specific spatial cluster typologies: High-High (hotspots), Low-Low (coldspots), and spatial anomalies (High-Low / Low-High).
\end{enumerate}

\subsection{Protective Action Decision Model (PADM) and Structural Equation Modeling (SEM)}
\label{subsec:padm_sem}

The societal value of meteorological forecasting depends on how effectively warning information guides public protective behavior. The Protective Action Decision Model (PADM) formulated by Lindell and Perry (2012)~\cite{lindell2012protective} provides a behavioral science framework detailing how environmental hazard cues propagate through risk perception, mediated by institutional trust, to drive protective action compliance.

In empirical behavioral modeling, Structural Equation Modeling (SEM) provides a simultaneous framework for validating latent measurement models and testing directional structural path hypotheses:
\begin{equation}
\label{eq:sem_measurement}
\mathbf{x} = \mathbf{\Lambda}_x \boldsymbol{\xi} + \boldsymbol{\delta}, \quad \mathbf{y} = \mathbf{\Lambda}_y \boldsymbol{\eta} + \boldsymbol{\epsilon}
\end{equation}
\begin{equation}
\label{eq:sem_structural}
\boldsymbol{\eta} = \mathbf{B} \boldsymbol{\eta} + \mathbf{\Gamma} \boldsymbol{\xi} + \boldsymbol{\zeta}
\end{equation}
where $\boldsymbol{\xi}$ and $\boldsymbol{\eta}$ represent exogenous and endogenous latent constructs, $\mathbf{x}$ and $\mathbf{y}$ are observed survey indicator vectors, $\mathbf{\Lambda}_x, \mathbf{\Lambda}_y$ are factor loading matrices, $\mathbf{B}, \mathbf{\Gamma}$ denote structural path coefficients, and $\boldsymbol{\delta}, \boldsymbol{\epsilon}, \boldsymbol{\zeta}$ are residual variance terms.

By administering an empirical field survey ($N=842$ across northern Chinese cities) and performing Harman's single-factor test to control for Common Method Bias (CMB), SEM quantitatively reveals how forecast certainty and institutional trust modulate citizen protective actions, establishing the foundation for calibrated municipal emergency response matrices.

---

\section{Critical Literature Review and Thesis Entry Points}
\label{sec:literature_gaps_synthesis}

\subsection{Comparative Analysis of Forecasting Paradigms}
\label{subsec:paradigms_comparison}

Reviewing decades of atmospheric dust forecasting evolution, Table~\ref{tab:paradigms_synthesis} summarizes four generations of technological paradigms alongside the proposed DustML framework.

\begin{table}[H]
\centering
\small
\caption{Comparative Matrix of Four Dust Forecasting Paradigms vs. Proposed DustML}
\label{tab:paradigms_synthesis}
\begin{tabularx}{\textwidth}{l p{2.2cm} p{2.2cm} p{2.3cm} p{2.2cm} X}
\toprule
\textbf{Dimension} & \textbf{Gen 1: Statistical} & \textbf{Gen 2: Dynamic NWP} & \textbf{Gen 3: Pure AI} & \textbf{Gen 4: Hybrid AI} & \textbf{This Work: DustML} \\
\midrule
Primary Models & Linear regression, ARIMA & ECMWF IFS, WRF-Chem, CMA-GFS & Shallow SVM, Random Forest, LSTM & ConvLSTM, PredRNN, unconstrained PINN & \textbf{Dual-Line: LightGBM Residual + 14-Node ST-GNN + PINN} \\
Input Sources & Single station surface gauges & Coarse grid atmospheric dynamics & Satellite pixels or surface telemetry & Grid-based radar or optical sequences & \textbf{Multi-Modal: NWP + ERA5 + FY-4A/MODIS on 5\,km grid} \\
Effective Lead Time & $< 24$ hours & $24\text{--}72$ hours (degrades $>72$h) & $< 48$ hours & $48\text{--}72$ hours & \textbf{Extended to 10--15 days (extended-range skill)} \\
Physical Consistency & Black-box empirical fitting & Follows PDEs; static empirical schemes & Poor (frequent calm-wind hallucinations) & Lacks conservation constraints & \textbf{Strict Owen saltation threshold and advective mass PDE ($>98\%$)} \\
Orographic Funneling & Cannot model & Smoothed topography blunts venturi jets & Euclidean 2D grids incur high overhead & Dense convolutions are computationally costly & \textbf{14-node bidirectional transport graph diffusion (8$\times$ faster)} \\
Societal Integration & None & Qualitative synoptic alerts & Deterministic points without confidence & Limited sociological integration & \textbf{Coupled GIS Kriging + PADM SEM + 4-tier emergency playbook} \\
\bottomrule
\end{tabularx}
\end{table}

\subsection{The Three Critical Research Gaps in the Literature}
\label{subsec:three_gaps}

Despite continuous advances, existing research across atmospheric sciences and environmental informatics leaves three critical scientific gaps unaddressed:

\begin{enumerate}
    \item \textbf{Gap 1: Orographic Topographic Smoothing and Systematic NWP Residuals} \\
    Global and regional NWP models ($0.1^\circ\text{--}0.25^\circ$) mechanically filter complex topography. Narrow mountain gaps (e.g., Hexi Corridor, Helan Mountain pass) serve as critical venturi channels for low-level dust jets. Coarse-grid smoothing severely dampens localized wind shear and saltation fluxes, causing systematic phase delays and peak-concentration underestimation in downstream metropolitan areas. Current literature lacks non-linear machine learning architectures capable of learning these systematic spatial residuals at extended-range lead times.

    \item \textbf{Gap 2: Unconstrained AI Physical Hallucinations and Conservation Violations} \\
    Contemporary deep learning architectures (e.g., standard CNNs, ConvLSTMs, Transformers) optimize empirical loss functions without physical boundary constraints. Because calm, clear-sky conditions dominate historical observational datasets, unconstrained networks memorize spurious correlations. Consequently, when evaluated on extreme out-of-distribution events, they frequently predict dust generation in calm conditions or create particulate mass from nothing, violating aerodynamic thresholding and mass conservation laws.

    \item \textbf{Gap 3: Institutional Disconnect Between Meteorological Predictions and Societal Response} \\
    Atmospheric engineering historically focuses exclusively on optimizing statistical validation metrics (such as RMSE), overlooking the operational translation of forecasts into public disaster mitigation. Current operational forecasts often provide vague deterministic contours without calibrated uncertainty bounds. Furthermore, the behavioral pathways linking physical hazard cues to citizen protective compliance remain largely unquantified, resulting in delayed public response and suboptimal municipal resource allocation during severe dust emergencies.
\end{enumerate}

\subsection{Methodological Breakthroughs of this Dissertation}
\label{subsec:thesis_breakthrough}

To address these three gaps, this dissertation implements three targeted methodological solutions:
\begin{itemize}
    \item \textbf{Resolving Gap 1}: Establishes Main Line A (tree-based residual learning via LightGBM and CatBoost coupled with asymmetric quantile regression), extracting non-linear topographic error corrections from multi-source reanalysis and surface observations to eliminate systematic NWP bias and generate $[P_{10}, P_{90}]$ uncertainty envelopes;
    \item \textbf{Resolving Gap 2}: Establishes Main Line B (an end-to-end deep spatiotemporal architecture integrating AI-GAMFS foundation priors, a 14-node corridor ST-GNN, and a continuous PINN regularizer), embedding Owen's (1964) saltation threshold gate and 2D advective mass continuity PDEs into the backpropagation graph to mathematically suppress physical hallucinations;
    \item \textbf{Resolving Gap 3}: Establishes an interdisciplinary disaster management framework that links calibrated forecast intervals to Ordinary Kriging spatial exposure mapping, LISA hotspot detection, and PADM structural equation modeling ($N=842$), formulating an operational 4-tier municipal emergency response decision matrix.
\end{itemize}

The theoretical equations, literature taxonomy, and scientific gaps detailed in this chapter establish the intellectual foundation for the data architectures and predictive models developed in subsequent chapters.
